% manual.tex -- TFR/911H NitrOS-9 Reference Pages.
%
% Set in the manner of Digital Equipment Corporation's ULTRIX reference
% pages of the early 1990s: Times text, Helvetica headings hung in the
% left margin, name(section) running heads.
%
% Build:  make -C .. manual     (or: xelatex manual.tex, twice)

\documentclass[10pt,twoside,openany]{book}
\usepackage[letterpaper,top=0.85in,bottom=0.9in,inner=1.05in,outer=0.95in,
            headheight=14pt,headsep=0.3in,footskip=0.45in]{geometry}
\usepackage{fontspec}
\setmainfont{Times}
\setsansfont{Helvetica}[Scale=0.92]
\setmonofont{Courier}[Scale=0.95]
\usepackage{fancyhdr}
\usepackage{xcolor}
\usepackage{array,tabularx,longtable,booktabs}
\usepackage{alltt}
\usepackage[hidelinks,pdftitle={TFR/911H NitrOS-9 Reference Pages},
            pdfauthor={The TFR/9 Project}]{hyperref}

\setlength{\parindent}{0pt}
\setlength{\parskip}{0.7ex plus 0.2ex}
\renewcommand{\arraystretch}{1.15}
\raggedbottom
\setlength{\emergencystretch}{2em}

% ── Little helpers ─────────────────────────────────────────────────
\newcommand{\cmd}[1]{\textbf{#1}}            % a command name
\newcommand{\var}[1]{\textit{#1}}            % something you supply
\newcommand{\file}[1]{\texttt{#1}}           % a path, address, register
\newcommand{\key}[1]{\textsf{\small #1}}     % a key: \key{Ctrl/C}
\newcommand{\rp}[2]{\textbf{#1}(#2)}         % a reference page
\newcommand{\type}[1]{\texttt{\textbf{#1}}}  % what you type, in text
\newcommand{\sans}[1]{\textsf{#1}}

% Examples: alltt, small Courier; \textbf for what you type.
\newenvironment{ex}{\par\vspace{0.3ex}\small\begin{alltt}}{\end{alltt}\vspace{0.3ex}}
\newenvironment{exsmall}{\par\vspace{0.3ex}\scriptsize\begin{alltt}}{\end{alltt}\vspace{0.3ex}}

% ── Running heads and feet ─────────────────────────────────────────
\newcommand{\thesectname}{}
\newcommand{\therefname}{}
\pagestyle{fancy}
\fancyhf{}
\renewcommand{\headrulewidth}{0pt}
\fancyhead[LE,RO]{\sffamily\bfseries\small\therefname}
\fancyhead[RE,LO]{\sffamily\bfseries\small\therefname}
\fancyfoot[LE,RO]{\sffamily\small\thepage}
\fancyfoot[RE,LO]{\sffamily\small\thesectname}
\fancypagestyle{plain}{\fancyhf{}\fancyfoot[LE,RO]{\sffamily\small\thepage}
  \fancyfoot[RE,LO]{\sffamily\small\thesectname}}

% Pages numbered chapter-page, as DEC numbered them: 1-1, 1-2, ... A-1.
\makeatletter
\newcommand{\chapterpages}{%
  \renewcommand{\thepage}{\thechapter--\arabic{page}}\setcounter{page}{1}}

% Chapter openings: a rule, the number big at the right, the title.
\renewcommand{\@makechapterhead}[1]{%
  \vspace*{0.4in}%
  {\color{black!75}\hrule height 2pt}\vspace{6pt}%
  \noindent\begin{minipage}[b]{0.78\textwidth}\raggedright
    \sffamily\bfseries\fontsize{24}{28}\selectfont #1\end{minipage}%
  \hfill{\sffamily\bfseries\fontsize{48}{48}\selectfont\color{black!70}\thechapter}\par
  \vspace{4pt}\hrule height 0.5pt\vspace{0.45in}}
\renewcommand{\@makeschapterhead}[1]{%
  \vspace*{0.4in}{\color{black!75}\hrule height 2pt}\vspace{6pt}%
  \noindent{\sffamily\bfseries\fontsize{24}{28}\selectfont #1}\par
  \vspace{4pt}\hrule height 0.5pt\vspace{0.45in}}

% Headings within chapters: Helvetica, numbered like 1.2.
\renewcommand{\section}{\@startsection{section}{1}{0pt}{3.2ex plus 1ex}{1.2ex}%
  {\sffamily\bfseries\large}}
\renewcommand{\subsection}{\@startsection{subsection}{2}{0pt}{2.4ex plus 0.6ex}{0.8ex}%
  {\sffamily\bfseries\normalsize}}
\setcounter{secnumdepth}{2}

% The table of contents, in Helvetica.
\renewcommand{\@pnumwidth}{2.8em}
\renewcommand{\@tocrmarg}{3.6em}
\renewcommand*{\l@chapter}[2]{\addpenalty{-\@highpenalty}\vskip 1.1em plus 1pt
  {\setlength{\@tempdima}{2em}\parindent\z@ \rightskip\@pnumwidth
   \parfillskip -\@pnumwidth \leavevmode\sffamily\bfseries
   \advance\leftskip\@tempdima \hskip -\leftskip #1\nobreak\hfil
   \nobreak\hb@xt@\@pnumwidth{\hss #2}\par}\vskip 0.3em}
\renewcommand{\l@section}{\@dottedtocline{1}{1.5em}{2.6em}}
\renewcommand{\l@subsection}{\@dottedtocline{2}{4.1em}{3.2em}}
\makeatother

% ── Reference pages ────────────────────────────────────────────────
% \refsection{4}{Special Files} starts a group of reference pages.
\newcommand{\refsection}[2]{%
  \clearpage
  \renewcommand{\therefname}{}%
  \renewcommand{\thesectname}{#2}%
  \section*{Section #1: #2}%
  \addcontentsline{toc}{section}{Section #1: #2}}

% \refentry{vi}{1}{a small vi} starts one reference page.
\newcommand{\refentry}[3]{%
  \clearpage
  \renewcommand{\therefname}{#1(#2)}%
  \phantomsection
  \addcontentsline{toc}{subsection}{#1(#2)\quad\textit{#3}}}

% The body of a reference page: headings hang in the left margin.
\newenvironment{refpage}{%
  \begin{list}{}{%
    \setlength{\leftmargin}{1.3in}\setlength{\labelwidth}{1.2in}%
    \setlength{\labelsep}{0.1in}\setlength{\itemsep}{1.6ex}%
    \setlength{\parsep}{0.7ex}\setlength{\topsep}{0pt}\setlength{\listparindent}{0pt}%
    \renewcommand{\makelabel}[1]{\sffamily\bfseries ##1\hfil}}}%
  {\end{list}}
\newcommand{\rh}[1]{\item[#1]}

% Two-column tables of keys or options inside a reference page.
\newenvironment{keys}[1][1.1in]{\par\tabularx{\linewidth}[t]{@{}>{\raggedright\arraybackslash}p{#1}>{\raggedright\arraybackslash}X@{}}}%
  {\endtabularx\par}

\begin{document}
\frontmatter
\pagestyle{empty}

% ── Title page ─────────────────────────────────────────────────────
\begin{titlepage}
\sffamily
\vspace*{0.9in}
{\fontsize{30}{36}\selectfont\bfseries TFR/911H NitrOS-9\\[4pt] Reference Pages\par}
\vspace{10pt}
{\Large Commands, Special Files, Games, and Maintenance\par}
\vspace{0.5in}
{\color{black!75}\hrule height 2pt}
\vspace{0.3in}
\begin{tabular}{@{}l@{}}
\bfseries Order Number: TFR-911H-RP-A
\end{tabular}

\vspace{0.3in}
\begin{minipage}{0.72\textwidth}
\rmfamily\normalsize
This manual describes the commands, games, and tools added to NitrOS-9
Level 2 on the TFR/911H single board computer, the hardware the TFR/911H
firmware presents to OS-9, and how to build and maintain the system. It
begins with an introduction to OS-9 Level 1 and Level 2 as they run on the
board.
\end{minipage}

\vspace{0.35in}
\begin{tabular}{@{}l l@{}}
\bfseries Revision/Update Information: & This is a new manual.\\[3pt]
\bfseries Operating System and Version: & NitrOS-9/6809 Level 2 V3.3.0\\[3pt]
\bfseries Hardware and Firmware: & TFR/911H, firmware v4\\
\end{tabular}

\vfill
{\bfseries\large the tfr/9 project\par}
{\small github.com/strickyak/tfr9\par}
\end{titlepage}

% ── Copyright page ─────────────────────────────────────────────────
\clearpage
\vspace*{\fill}
{\small
\textbf{First Printing, October 2026}

The information in this document is subject to change without notice and
should not be construed as a commitment by anyone in particular. The
software described in this document is furnished as source code, and may
be used, copied, and changed under the terms of the licenses that come
with it.

No responsibility is assumed for the use or reliability of software on
equipment that was soldered together at a kitchen table.

\copyright{} The TFR/9 Project contributors, 2026. All rights reserved.

\medskip
\textbf{Trademarks.} OS-9 is a trademark of Microware Systems Corporation.
NitrOS-9 is the work of the NitrOS-9 Project. Zork is a registered
trademark of Infocom, Inc. Tetris is a trademark of The Tetris Company.
Raspberry Pi is a trademark of Raspberry Pi Ltd. Color Computer and
CoCo are associated with Tandy Corporation. ULTRIX and DEC are
trademarks of Digital Equipment Corporation.

\medskip
\textbf{About the design.} This manual is typeset in the manner of
Digital Equipment Corporation's ULTRIX reference pages of the early 1990s,
as a tribute. It is not a Digital document, and it has nothing to do with
Digital Equipment Corporation or its successors.

\medskip
This document was prepared with \LaTeX{} from
\file{v4/nitros9-l2/doc/manual.tex}.
}
\clearpage

\pagestyle{fancy}
\renewcommand{\thesectname}{Contents}
\setcounter{page}{3}
\tableofcontents

% ── Preface ────────────────────────────────────────────────────────
\chapter*{About This Manual}
\addcontentsline{toc}{chapter}{About This Manual}
\renewcommand{\thesectname}{About This Manual}

This manual describes the software that turns a TFR/911H into a small,
complete OS-9 system: NitrOS-9 Level 2, booted by the board's firmware,
with a disk, a clock, a Unix-flavored set of commands, an editor, a sound
chip, and games. It also describes the hardware that the firmware
emulates for OS-9, and how to rebuild everything.

\section*{Audience}
This manual is for anyone with a TFR/911H on their desk: people who know
some Unix and want to find their way around OS-9, CoCo veterans who want
to know what is different about this machine, and people who want to
change the firmware or the disk. No knowledge of OS-9 is assumed in
Chapter~1.

\section*{Organization}
This manual has two chapters and two appendices:

\begin{tabularx}{\linewidth}{@{}>{\bfseries}l X@{}}
Chapter 1 & Introduces the TFR/911H, OS-9 Level 1 and Level 2, the
TurbOS and NitrOS-9 systems that run on the board, and a first
session.\\
Chapter 2 & Contains the reference pages, grouped as ULTRIX grouped
them: Section 1, Commands; Section 4, Special Files (the hardware);
Section 6, Games; and Section 8, Maintenance.\\
Appendix A & Explains how to connect a speaker.\\
Appendix B & Describes how the software was tested.\\
\end{tabularx}

\section*{Related Documents}
\begin{itemize}\setlength{\itemsep}{0pt}
\item \file{v4/nitros9-l2/README.md}, which this manual expands on.
\item The NitrOS-9 documentation, at \file{nitros9.org} and in the
NitrOS-9 source tree under \file{docs}.
\item \textit{OS-9 Level Two} for the Tandy Color Computer 3, whose
commands NitrOS-9 still follows closely.
\item The TFR/9 hardware reference, on the \file{strickyak/tfr9} wiki.
\end{itemize}

\section*{Conventions}
\begin{tabularx}{\linewidth}{@{}>{\raggedright\arraybackslash}p{1.3in}X@{}}
\toprule
\textbf{Convention} & \textbf{Meaning}\\
\midrule
\texttt{OS9:} & The NitrOS-9 Shell prompt.\\
\texttt{TCL>} & The TFR/911H firmware's console prompt, before an OS boots.\\
\texttt{\$} & A shell prompt on the computer the board is plugged into.\\
\type{bold type} & In examples, what you type.\\
\var{italic} & Something you supply: a file name, a number.\\
\texttt{[ ]} & In syntax, an optional part. Do not type the brackets.\\
\texttt{|} & In syntax, a choice: \texttt{on|off} means \texttt{on} or \texttt{off}.\\
\key{Ctrl/C} & Hold down the Ctrl key while you press C. Also written \texttt{\^{}C}.\\
\key{Return} & The Return (Enter) key.\\
\rp{cmd}{n} & The reference page for \cmd{cmd} in Section~\var{n}.\\
\file{\$FF41} & A hexadecimal number, often an address in the 6309's I/O page.\\
\bottomrule
\end{tabularx}

% ═══════════════════════════════════════════════════════════════════
\mainmatter
\chapter{Introduction}
\chapterpages
\renewcommand{\thesectname}{Introduction}

This chapter introduces the TFR/911H, the two levels of OS-9, the two
OS-9 systems the board runs, and what was added to NitrOS-9 Level 2 for
it. It ends with a first session at the keyboard.

\section{The TFR/911H}

The TFR/911H is a single board computer built around a genuine Hitachi
HD6309 microprocessor---the CMOS cousin of Motorola's 6809, with extra
registers and a faster native mode---and a Raspberry Pi RP2350B
microcontroller. Everything that is not the CPU is done by the RP2350B:
\begin{itemize}\setlength{\itemsep}{0pt}
\item It makes the 6309's E and Q clocks, one bus cycle at a time.
\item It is the memory: 256K of the RP2350B's 520K of SRAM, behind an
emulated MMU.
\item It is every I/O device: a serial port for the console, a hard disk,
a clock, interrupt timers, a sound chip.
\end{itemize}
One of the RP2350B's two cores runs the 6309's bus; the other talks USB
to a program on your computer, the \emph{tether}, which is the board's
terminal, its disk drive, and its clock. The 6309 runs at about 2.8 to
3.2~MHz---faster than the 1.79~MHz of the fastest Color Computer.

The TFR/9 family of boards was designed to make very cheap 6809-family
computers that run OS-9. Its design files and software are at
\file{github.com/strickyak/tfr9}.

\section{OS-9 Level 1 and Level 2}

OS-9 was written by Microware Systems Corporation in 1979 and 1980 for
the Motorola 6809. It is a real-time, multitasking, multi-user operating
system that will look familiar to anyone who knows Unix: a hierarchical
file system, a shell with pipes and redirection, devices that are files,
processes with parents and children. It was written to be small and to
run from ROM, so everything in OS-9 is a \emph{module}: a
position-independent, re-entrant block of code or data with a header and
a CRC, which the system can find in memory by name. The kernel, the
device drivers, and the commands you run are all modules.

The 6809 can address only 64K of memory, and OS-9 comes in two
\emph{levels} that differ in what they do about that.

\begin{description}\setlength{\itemsep}{0.5ex}
\item[\sans{Level 1}] runs everything in a single 64K address space: the
kernel, every process, every module. It needs no memory management
hardware, which made it right for the first Color Computers. Tandy sold
OS-9 Level One for the Color Computer in 1983.

\item[\sans{Level 2}] uses a memory management unit to give each process
its own 64K address space, made of 8K blocks drawn from a larger physical
memory---up to 2M. The kernel lives in a space of its own. Processes cannot
see or damage each other's memory, and many more can run at once. Tandy
sold OS-9 Level Two for the Color Computer~3, whose GIME chip has the
MMU, in 1987.
\end{description}

NitrOS-9 began in the 1990s as a rewrite of OS-9 that made use of the
Hitachi 6309's extra registers and instructions. Today the NitrOS-9
Project maintains it as open source, for both the 6309 and the 6809, for
Level 1 and Level 2, and for many machines besides the Color Computer.

\section{Level 1 on the TFR/911H: TurbOS}

When the TFR/911H firmware starts, it waits at its console prompt,
\texttt{TCL>}, a small Tcl interpreter (see \rp{tcl}{8}). Typing
\type{bye} there starts \emph{TurbOS}, a small OS-9 Level 1 system that
the firmware loads into the 6309's memory as a single 64K image, from its
own flash.

TurbOS runs without a disk. Its kernel, I/O manager, SCF terminal manager,
terminal driver, Shell, and a handful of commands (\cmd{mdir}, \cmd{mfree},
\cmd{procs}, \cmd{sleep}) are modules in the image. Its prompt is
\texttt{TOS:}:
\begin{ex}
TOS:\textbf{mdir}
 Module directory at 00:00:59
kernel  init  tk  ioman  scf  scvt  term  go  shell  mfree  mdir  procs  sleep
\end{ex}
TurbOS comes from the TFR/9 project's TurboLab work, in the
\file{v4/turbolab} directory of the TFR/9 source. The rest of this manual is about Level 2.

\section{Level 2 on the TFR/911H: NitrOS-9 V3.3.0}

Typing \type{l2} at the \texttt{TCL>} prompt starts NitrOS-9/6809 Level 2
V3.3.0, built for the Color Computer~3. The firmware copies a small
launcher and the kernel's boot track into memory and resets the 6309,
which boots in about half a second:
\begin{ex}
NitrOS-9/6809 Level 2 V3.3.0
Tandy Color Computer 3
(C) 2014 The NitrOS-9 Project
Fri Oct 16 17:23:34 2026
http://www.nitros9.org

NitrOS-9 Level 2 on the TFR/911H.  Unix-style commands: man unix

Shell

OS9:
\end{ex}

\subsection{The machine NitrOS-9 sees}

NitrOS-9 thinks it is on a Color Computer~3. The firmware emulates as
much of one as OS-9 uses, and a little more:

\begin{tabularx}{\linewidth}{@{}l l X@{}}
\toprule
\textbf{Device} & \textbf{Address} & \textbf{What it is}\\
\midrule
RAM & & 256K in 32 blocks of 8K; 184K free for programs\\
GIME MMU & \file{\$FF90}--\file{\$FF91}, \file{\$FFA0}--\file{\$FFAF} & Two task maps; the constant page at \file{\$FE00}\\
GIME timer & \file{\$FF92} & The VBORD interrupt, 60 times a second: the system clock\\
6850 ACIA & \file{\$FF14}--\file{\$FF15} & The console, \file{/Term}\\
SN76489 & \file{\$FF41} & A sound chip: three tones and noise\\
Control & \file{\$FF42} & Keyclicks, the bell, raw arrow keys\\
Info port & \file{\$FF43}--\file{\$FF44} & The firmware's report on itself\\
Becker port & \file{\$FF45}--\file{\$FF46} & DriveWire to a FujiNet: disks \file{/X0}--\file{/X3}, network \file{/N}\\
EmuDsk & \file{\$FF80}--\file{\$FF86} & A virtual hard disk, \file{/DD} and \file{/H0}\\
Clock & \file{\$FFC0}--\file{\$FFC7} & Date and time, from your computer\\
\bottomrule
\end{tabularx}

\rp{tfr911h}{4} describes each of these.

\subsection{The disk}

The disk is a file on your computer, \file{v4/nitros9-l2/level2.dsk}: an
RBF disk of 99,999 sectors (24.4M) that the tether serves to EmuDsk a
sector at a time. It holds \file{OS9Boot}, the rest of the operating
system, which the kernel loads at boot; the commands in \file{/DD/CMDS};
the system files \file{errmsg} and \file{helpmsg} in \file{/DD/SYS}; and
story files for \rp{zm}{6} in \file{/DD/GAMES}. Whatever you save goes to
the file at once.

With a FujiNet attached, typing \type{l2fn} instead of \type{l2} boots from
the FujiNet's first device slot, a copy of this disk on its SD card or a
TNFS server; \file{/DD} is then the DriveWire disk \file{/X0}, and the
tether's disk is \file{/H0}. See \rp{fujinet}{4}.

\subsection{The modules}

These are in memory once the system is up (\rp{info}{1} lists them):

\begin{tabularx}{\linewidth}{@{}>{\raggedright\arraybackslash}p{1.45in}X@{}}
\toprule
\textbf{Modules} & \textbf{What they do}\\
\midrule
\file{REL Boot Krn} & Relocate the boot track, load \file{OS9Boot} from EmuDsk (or, for \cmd{l2fn}, the FujiNet), and the kernel\\
\file{Init} & The system's settings: startup program, devices, table sizes\\
\file{KrnP2 KrnP3} & The rest of the kernel, and its error-message extension\\
\file{IOMan} & The I/O manager\\
\file{Clock Clock2} & The 60~Hz tick, and the real-time clock (\file{clock2\_tfr911})\\
\file{SCF sc6850 Term} & The terminal file manager, the 6850 driver, \file{/Term}\\
\file{RBF EmuDsk DD H0} & The disk file manager, the EmuDsk driver, \file{/DD} and \file{/H0}\\
\file{rbdw dwio X0}--\file{X3} & DriveWire disks on the FujiNet, and DriveWire's wire routines\\
\file{scdwv N N1}--\file{N5} & DriveWire's network devices, through which the \cmd{fn} commands talk\\
\file{PipeMan Piper Pipe} & Pipes\\
\file{SysGo Shell} & The first process, and the Shell it starts\\
\bottomrule
\end{tabularx}

\section{What Was Added}

These are the commands, games, and tools described in Chapter~2. Each has
a \cmd{man} page on the board too: type \type{man} \var{name}.

\begin{tabularx}{\linewidth}{@{}>{\raggedright\arraybackslash}p{1.25in}X@{}}
\toprule
\textbf{Reference page} & \textbf{Summary}\\
\midrule
\rp{unix}{1} & \cmd{ls}, \cmd{cd}, \cmd{cat}, \cmd{cp}, \cmd{mv}, \cmd{rm},
\cmd{ps}, \cmd{df}, and other Unix names for OS-9 commands\\
\rp{vi}{1} & A small vi, also called \cmd{vim}\\
\rp{info}{1} & Everything the system knows about itself\\
\rp{grep}{1} & Print lines that contain a string\\
\rp{pwd}{1} & Print the data or execution directory\\
\rp{clear}{1} & Clear the screen\\
\rp{beep}{1}, \rp{click}{1} & Ring the bell; turn keyclicks on and off\\
\rp{fujinet}{1} & The FujiNet's hosts, device slots, and status: \cmd{fnstatus}, \cmd{fnlistdevs}, \ldots\\
\rp{fnconfig}{1} & FujiNet CONFIG\\
\rp{fnweather}{1} & The weather where you are\\
\rp{fujinet}{4} & The FujiNet: its wiring, its disks, and booting from it\\
\rp{tfr911h}{4} & The hardware the firmware emulates\\
\rp{sn76489}{4} & Programming the sound chip\\
\rp{tetris}{6} & Tetris, with music\\
\rp{zm}{6} & A Z-machine, for Zork and the other Infocom games\\
\rp{tether}{8} & The program that is the board's terminal and disk\\
\rp{tcl}{8} & The firmware's console: \cmd{l2}, \cmd{l2fn}, \cmd{bye}, \cmd{fujinet}, \cmd{beep}, \cmd{click}\\
\rp{nitros9-l2}{8} & Building the system and keeping a disk up to date\\
\bottomrule
\end{tabularx}

\section{A First Session}

Plug the TFR/911H into a USB port and, in the \file{v4} directory of the
TFR/9 source, run \file{run-l2.sh}. It starts the tether, which boots
NitrOS-9 by typing \type{l2} for you (or reconnects, if NitrOS-9 is
already running):
\begin{ex}
$ \textbf{WIRE=/dev/cu.usbmodem31201 sh run-l2.sh}
TFR/911H on /dev/cu.usbmodem31201, disk .../nitros9-l2/level2.dsk.
^C goes to OS-9; quit with ^] .
\end{ex}
Then look around. \cmd{man unix} lists the Unix-style names; \cmd{ls -l}
lists a directory; \cmd{info} shows the whole system:
\begin{ex}
OS9:\textbf{ls -l /dd/cmds}
OS9:\textbf{cd /dd}
OS9:\textbf{pwd}
/DD
OS9:\textbf{info}
\end{ex}
Write a file with \cmd{vi}, and look at it with \cmd{cat}:
\begin{ex}
OS9:\textbf{vi notes}
OS9:\textbf{cat notes}
\end{ex}
Play something:
\begin{ex}
OS9:\textbf{zm zork1}
OS9:\textbf{tetris}
\end{ex}
\key{Ctrl/C} interrupts the program that is running, as on Unix. To leave
the tether, press \key{Ctrl/]}; NitrOS-9 keeps running on the board, and
\file{run-l2.sh} will reconnect to it.

OS-9 differs from Unix in a few ways you will meet early:
\begin{itemize}\setlength{\itemsep}{0pt}
\item Each process has two current directories: the \emph{data}
directory, for files (\cmd{cd}, \cmd{pwd}), and the \emph{execution}
directory, where commands are found (\cmd{chx}, \cmd{pxd}). Both start
at \file{/DD} and \file{/DD/CMDS}.
\item The pipe symbol is \texttt{!}, not \texttt{|}. \texttt{>>} redirects
standard error; there is no append. \texttt{\&} runs a command in the
background, as on Unix.
\item File names are not case sensitive. Text files end their lines with
a carriage return, not a line feed.
\end{itemize}

% ═══════════════════════════════════════════════════════════════════
\chapter{Reference Pages}
\chapterpages
\renewcommand{\thesectname}{Reference Pages}

The reference pages follow, in four sections. Within a section they are in
alphabetical order. Each page has the same parts, where they apply:
\textbf{Name}, \textbf{Syntax}, \textbf{Description}, \textbf{Options},
\textbf{Restrictions}, \textbf{Examples}, \textbf{Files}, and \textbf{See
Also}. The commands that came with NitrOS-9 are not described here; use
\cmd{man} on the board for those.

% ─────────────────────────────────────────────────────────────────── 1
\refsection{1}{Commands}
This section describes the commands added to NitrOS-9 Level~2 for the TFR/911H. The commands that come with NitrOS-9 are described by \cmd{man} on the board, and in the NitrOS-9 documentation.

\begin{tabularx}{\linewidth}{@{}>{\bfseries}p{1.3in}X@{}}
\toprule
\textnormal{\textbf{Name}} & \textbf{Description}\\
\midrule
beep(1) & Ring the bell\\
clear(1) & Clear the screen\\
click(1) & Turn keyclicks on or off\\
fnconfig(1) & FujiNet CONFIG: the FujiNet's status, hosts, device slots, and WiFi\\
fnweather(1) & The weather where you are, through the FujiNet\\
fujinet(1) & Commands for the FujiNet: \textmd{fnstatus, fnlisthosts, fnlistdevs, fnsetdevfile,} \ldots\\
grep(1) & Print lines that contain a string\\
info(1) & Show what the system and the hardware know about themselves\\
pwd(1) & Print the data or execution directory\\
unix(1) & Unix-style names for OS-9 commands: \textmd{ls, cd, cat, cp, mv, rm, ps, df,} \ldots\\
vi(1) & A small vi, also called \textmd{vim}\\
\bottomrule
\end{tabularx}


\refentry{beep}{1}{ring the bell}
\begin{refpage}
\rh{Name} \cmd{beep} -- ring the bell
\rh{Syntax} \cmd{beep}
\rh{Description}
The \cmd{beep} command writes a bell character (\file{\$07}) to standard
output. The terminal the tether runs in rings its bell, and the TFR/911H
firmware, which sees the character go out through the 6850, plays an
880~Hz tone for 150~ms through the board's speaker, if one is connected
and bit~1 of the control register at \file{\$FF42} is set (it is, at
power on).

Any program that prints \texttt{\^{}G} makes the same sound.
\rh{Examples}
To sound the bell when a long command finishes:
\begin{ex}
OS9:\textbf{copy /dd/bigfile /dd/backup; beep}
\end{ex}
\rh{See Also} \rp{click}{1}, \rp{sn76489}{4}, \rp{tfr911h}{4}
\end{refpage}

\refentry{clear}{1}{clear the screen}
\begin{refpage}
\rh{Name} \cmd{clear} -- clear the screen
\rh{Syntax} \cmd{clear}
\rh{Description}
The \cmd{clear} command sends the ANSI sequences that move the cursor home
and clear the screen of the terminal the tether runs in.
\rh{See Also} \rp{unix}{1}
\end{refpage}

\refentry{click}{1}{turn keyclicks on or off}
\begin{refpage}
\rh{Name} \cmd{click} -- turn keyclicks on or off
\rh{Syntax} \cmd{click} \texttt{on|off}
\rh{Description}
When keyclicks are on, the TFR/911H's speaker gives a short tick each time
a key reaches OS-9. The \cmd{click} command turns them on or off by setting
or clearing bit~0 of the control register at \file{\$FF42}. They are on at
power on. With no argument, or any other, \cmd{click} prints its usage.

The \texttt{TCL>} console has a \cmd{click} command of its own, which does
the same before an OS is booted; see \rp{tcl}{8}.
\rh{Restrictions}
\rp{tetris}{6} turns keyclicks off while it runs, and puts them back as
they were when it ends.
\rh{See Also} \rp{beep}{1}, \rp{sn76489}{4}, \rp{tfr911h}{4}
\end{refpage}

\refentry{fnconfig}{1}{FujiNet CONFIG}
\begin{refpage}
\rh{Name} \cmd{fnconfig} -- FujiNet CONFIG for NitrOS-9
\rh{Syntax} \cmd{fnconfig}
\rh{Description}
The \cmd{fnconfig} command is a full-screen program for managing the
FujiNet, as the CONFIG program does on other computers. It has four
pages, chosen with the number keys:

\begin{keys}[0.85in]
\key{1} Status & The FujiNet's firmware version, host name, WiFi network,
and IP settings.\\
\key{2} Hosts & The eight host slots. \key{Up} and \key{Down} pick one;
\key{E} edits it, \key{M} mounts it, and \key{U} unmounts it.\\
\key{3} Devices & The four DriveWire device slots, \file{/X0}--\file{/X3}.
\key{B} browses the slot's host for an image, \key{E} edits the path,
\key{M} mounts the image (it asks whether read-only or read/write), and
\key{U} ejects it.\\
\key{4} WiFi & \key{S} scans for networks; \key{E} enters a network's
name. Either then asks for the password.\\
\end{keys}

On every page \key{R} reads everything from the FujiNet again, \key{I}
shows the program's and the firmware's versions, \key{?} shows help, and
\key{Q} quits. The FujiNet saves each change when it succeeds.
\rh{Restrictions}
When NitrOS-9 was booted from the FujiNet (\cmd{l2fn}), \file{/X0} is
\file{/DD}: \cmd{fnconfig} will not change device slot~0 or its host, and
asks before changing the WiFi, which may drop the connection the disk
depends on.

The screen must be at least 80 columns by 24 lines. \cmd{fnconfig} was
written for a Color Computer~3's window; on the TFR/911H it is built to
send ANSI escape sequences instead.
\rh{See Also} \rp{fujinet}{1}, \rp{fujinet}{4}
\end{refpage}

\refentry{fnweather}{1}{the weather where you are}
\begin{refpage}
\rh{Name} \cmd{fnweather} -- the weather where you are, through the FujiNet
\rh{Syntax} \cmd{fnweather}
\rh{Description}
The \cmd{fnweather} command finds where you are from your network's public
address (from \texttt{ip-api.com}), then shows the weather there now (from
Open-Meteo): the sky, the temperature and what it feels like, the
humidity, the air pressure, and the wind. It uses Fahrenheit and miles in
the United States, Liberia, and Myanmar, and Celsius and kilometres
elsewhere. The FujiNet makes the HTTPS requests and picks the fields out
of the JSON replies.

\begin{keys}[0.5in]
\key{R} & Fetch the weather again.\\
\key{U} & Change between metric and imperial units.\\
\key{L} & Find where you are again.\\
\key{D} & Show the data as received, and the last error.\\
\key{I}, \key{?} & Show the credits; show help.\\
\key{Q} & Quit.\\
\end{keys}
\rh{Restrictions}
Each request can take several seconds. \cmd{fnweather} shows today's
conditions, not a forecast.
\rh{See Also} \rp{fnconfig}{1}, \rp{fujinet}{4}
\end{refpage}

\refentry{fujinet}{1}{commands for the FujiNet}
\begin{refpage}
\rh{Name} \cmd{fnstatus}, \cmd{fnlisthosts}, \cmd{fngethost},
\cmd{fnsethost}, \cmd{fnmount}, \cmd{fnlistdevs}, \cmd{fngetdevfile},
\cmd{fnsetdevfile}, \cmd{fnmountimg} -- manage the FujiNet
\rh{Syntax} \cmd{fnstatus}\\
\cmd{fnlisthosts}\\
\cmd{fngethost} \var{host\_slot}\\
\cmd{fnsethost} \var{host\_slot} \var{host}\\
\cmd{fnmount} \var{host\_slot}\\
\cmd{fnlistdevs}\\
\cmd{fngetdevfile} \var{device\_slot}\\
\cmd{fnsetdevfile} \var{device\_slot} \var{host\_slot} \var{mode} \var{path}\\
\cmd{fnmountimg} \var{device\_slot} [\var{mode}]
\rh{Description}
These commands talk to the FujiNet through \file{/N1}. A FujiNet has eight
\emph{host slots}, each a place images come from: a TNFS server, or
\texttt{SD} for its own SD card. It has four DriveWire \emph{device
slots}, each an image from a host, which NitrOS-9 sees as the disks
\file{/X0}--\file{/X3}. Slots are numbered from 0, as the FujiNet numbers
them: device slot~0 is \file{/X0}. A \var{mode} is 1 for read-only or 2
for read/write.

\begin{keys}[1.05in]
\cmd{fnstatus} & Print the FujiNet's firmware version, host name, WiFi
network, IP address, gateway, netmask, and name server, and whether the
WiFi is connected.\\
\cmd{fnlisthosts} & List the host slots.\\
\cmd{fngethost} & Print one host slot.\\
\cmd{fnsethost} & Set one host slot, keeping the others.\\
\cmd{fnmount} & Mount a host, so that its images can be used.\\
\cmd{fnlistdevs} & List the device slots: each one's host slot, mode, and
image.\\
\cmd{fngetdevfile} & Print one device slot's image.\\
\cmd{fnsetdevfile} & Put the image \var{path}, from host
\var{host\_slot}, in a device slot. The FujiNet saves it.\\
\cmd{fnmountimg} & Mount a device slot's image (read-only, unless
\var{mode} is 2).\\
\end{keys}
\rh{Examples}
Use an image on the FujiNet's SD card (host slot~1) as \file{/X1}:
\begin{ex}
OS9: \textbf{fnsetdevfile 1 1 2 /games/disk.dsk}
OS9: \textbf{fnmountimg 1 2}
OS9: \textbf{dir /x1}
\end{ex}
\rh{See Also} \rp{fnconfig}{1}, \rp{fujinet}{4}
\end{refpage}

\refentry{grep}{1}{print lines that contain a string}
\begin{refpage}
\rh{Name} \cmd{grep} -- print lines that contain a string
\rh{Syntax} \cmd{grep} [\texttt{-i}] [\texttt{-n}] \var{pattern} [\var{path}]
\rh{Description}
The \cmd{grep} command prints each line of \var{path}, or of standard input
if no path is given, that contains \var{pattern}. The pattern is a plain
string, not a regular expression. It may be given in double quotes, which
it needs if it contains spaces, or without them.

This is NitrOS-9's \cmd{grep}, changed in two ways: the pattern need not be
quoted, and it ends cleanly at the end of a file, where the original
reported a spurious error.
\rh{Options}
\begin{keys}[0.5in]
\texttt{-i} & Ignore case.\\
\texttt{-n} & Put each line's number in front of it.\\
\end{keys}
\rh{Examples}
\begin{ex}
OS9:\textbf{grep -n total sample.c}
OS9:\textbf{dir -e ! grep -i tetris}
\end{ex}
\rh{See Also} \rp{unix}{1}
\end{refpage}

\refentry{info}{1}{show the system and the hardware}
\begin{refpage}
\rh{Name} \cmd{info} -- show what the system and the hardware know about themselves
\rh{Syntax} \cmd{info} [\texttt{-p}]
\rh{Description}
The \cmd{info} command prints a report in eight parts:

\begin{keys}[0.85in]
\sans{System} & The OS name and version, the machine, the kernel's level
and version and the CPU it was built for, the startup program, disk and
console, the sizes of the device and interrupt tables, and the date. These
come from the \file{Init} module.\\
\sans{Memory} & The RAM, what is free (and the largest piece that is
free), and what is used by modules and by everything else, from the
system's memory block map.\\
\sans{Processes} & Every process: its number, its parent, its user, its
priority, its memory, its state, and the program it runs. The one marked
\texttt{*} is \cmd{info} itself.\\
\sans{Devices} & Every device descriptor in memory, with its file manager,
driver, port, and drive, and how many paths are open on it. A device that
nothing is using is not attached, and shows \texttt{-}.\\
\sans{Disks} & Each disk drive (once, though several descriptors may name
it): its name, the date it was made, its size, and its free space, from
its sector~0 and its allocation bitmap.\\
\sans{Interrupts} & The entries in the IRQ polling table, with the driver
and device each belongs to.\\
\sans{Modules} & The names of the modules in memory, and their total size.\\
\sans{Hardware} & The TFR/911H firmware's own report: the board, the
firmware's build date, the RP2350B's clock, memory and ID, the 6309's
measured speed and how busy it is, the interrupt rate, the disk, console,
sound, and FujiNet counters, whether the clock was set, and the uptime.\\
\end{keys}

The kernel's tables are in the system's own address space. \cmd{info}
reads them as NitrOS-9's \cmd{procs}, \cmd{devs}, and \cmd{mdir} do: it
gets the system process's descriptor with \texttt{F\$GPrDsc}, and uses
the memory map in it to copy what it needs with \texttt{F\$CpyMem}. The
hardware report comes from the firmware's info port; see \rp{tfr911h}{4}.
\rh{Options}
\begin{keys}[0.5in]
\texttt{-p} & Plain text: leave out the escape sequences that make the
headings bold, for saving the report to a file.\\
\end{keys}
\rh{Restrictions}
The RP2350B's temperature sensor reads nonsense on some boards, so the
chip's temperature is shown only when its reading is believable.

On firmware without the info port, the hardware part says so and is
otherwise empty.
\rh{Examples}
Part of a report, taken with \rp{tetris}{6} running in the background:
\begin{exsmall}
OS9:\textbf{info}
NitrOS-9/6809 Level 2 V03.3.0 on the TFR/911H

System
  OS          NitrOS-9/6809 Level 2 V03.3.0
  Machine     Tandy Color Computer 3, as the TFR/911H pretends to be
  Kernel      Level 2, version 3.3.0, built for the 6809
  Startup     runs SysGo, disk /DD, console /Term
  Tables      39 device slots, 21 IRQ polling slots
  Date        2026/10/17 01:25:15

Memory
  RAM         248K in 31 blocks of 8K
  Free        160K (largest piece 152K)
  In use      24K modules, 64K system and processes

Processes
   ID  Parent  User  Prio  Memory  State      Program
   1     0       0     255   0                  (the system)
   2     1       0     128   768     waiting    Shell
   3     2       0     128   4K      waiting    tetris
  *4     2       0     128   6K      running    info

Devices
  Device    Type  Manager  Driver    Port   Drive  Users
  /Term       SCF   SCF      sc6850    $FF14         4
  /DD         RBF   RBF      EmuDsk    $FF80  0      5
  /H0         RBF   RBF      EmuDsk    $FF80  0      -
  /Pipe       Pipe  PipeMan  Piper     $0000         -

Disks
  /DD = /H0
              "tfr911-level2.dsk", made 2026/10/17
              24.4M (99999 sectors), 24.0M free (98%)

Hardware (from the TFR/911H firmware)
  Board       TFR/911H: a Hitachi 6309 run by a Raspberry Pi RP2350B
  Pico        RP2350B at 250 MHz, 520K SRAM
  6309        3.24 MHz (measured), 70% busy, in 6809 mode
  Interrupts  59 IRQs/s: GIME VBORD 60 Hz clock, 6850 ACIA
  Sound       SN76489 at $FF41 (182 writes), PWM audio on GPIO 8, keyclick off
  Uptime      Pico 0m 56s, NitrOS-9 0m 54s
\end{exsmall}
To keep a report:
\begin{ex}
OS9:\textbf{info -p >/dd/sysinfo}
\end{ex}
\rh{See Also} \rp{tfr911h}{4}; NitrOS-9's \cmd{mdir}, \cmd{mfree}, \cmd{procs}, \cmd{devs}, \cmd{irqs}, \cmd{free}
\end{refpage}

\refentry{pwd}{1}{print the data or execution directory}
\begin{refpage}
\rh{Name} \cmd{pwd}, \cmd{pxd} -- print the data or execution directory
\rh{Syntax} \cmd{pwd}\\ \cmd{pxd}
\rh{Description}
An OS-9 process has two current directories. The \cmd{pwd} command prints
the \emph{data} directory, where file names that do not begin with
\texttt{/} are looked for, and which \cmd{cd} (or \cmd{chd}) changes. The
\cmd{pxd} command prints the \emph{execution} directory, where commands
are looked for, which \cmd{chx} changes.

Both are NitrOS-9's \cmd{pd}, built twice.
\rh{Examples}
\begin{ex}
OS9:\textbf{cd /dd/games}
OS9:\textbf{pwd}
/DD/GAMES
OS9:\textbf{pxd}
/DD/CMDS
\end{ex}
\rh{See Also} \rp{unix}{1}
\end{refpage}

\refentry{unix}{1}{Unix-style names for OS-9 commands}
\begin{refpage}
\rh{Name} \cmd{unix} -- Unix-style names for OS-9 commands
\rh{Syntax} \cmd{ls} [\texttt{-l}] [\var{path}]\quad \cmd{cd} \var{path}\quad \cmd{cat} \var{path}\quad \ldots
\rh{Description}
The TFR/911H's NitrOS-9 knows the names a Unix user reaches for first. Most
are \emph{aliases}: tiny modules that run the OS-9 command that does the
job, passing along their arguments. Two are built into the Shell, and the
rest are real commands. \type{man unix} on the board lists them.

\begin{tabularx}{\linewidth}{@{}l l X@{}}
\toprule
\textbf{Name} & \textbf{Runs} & \textbf{What it does}\\
\midrule
\cmd{ls} [\texttt{-l}] & \cmd{dir} [\texttt{-e}] & List a directory; \texttt{-l} gives the long form\\
\cmd{ll} & \cmd{dir -e} & List a directory, long form\\
\cmd{cat} & \cmd{list} & Print a text file\\
\cmd{cp} & \cmd{copy} & Copy a file\\
\cmd{mv} & \cmd{rename} & Rename a file or directory\\
\cmd{rm} & \cmd{del} & Delete files\\
\cmd{mkdir} & \cmd{makdir} & Make a directory\\
\cmd{rmdir} & \cmd{deldir} & Delete a directory and what is in it\\
\cmd{ps} & \cmd{procs} & List processes; \texttt{-e} lists everyone's\\
\cmd{df} & \cmd{free} & Show free space on a disk\\
\cmd{less} & \cmd{more} & Show a file a screen at a time\\
\cmd{man} & \cmd{help} & Show help for a command\\
\cmd{uname} & \cmd{echo} & Say what this is\\
\cmd{vim} & \cmd{vi} & The editor; see \rp{vi}{1}\\
\midrule
\cmd{cd} & built in & Change the data directory (the same as \cmd{chd})\\
\cmd{exit} & built in & Leave this Shell\\
\midrule
\cmd{pwd}, \cmd{pxd} & \cmd{pd} & See \rp{pwd}{1}\\
\cmd{grep}, \cmd{clear} & & See \rp{grep}{1}, \rp{clear}{1}\\
\bottomrule
\end{tabularx}

An alias puts any extra arguments of its own in front of yours, and then
chains to the real command with \texttt{F\$Chain}, so it costs no extra
process. \cmd{cd} and \cmd{exit} are added to the table of built-in
commands in a copy of the Shell's source.
\rh{Restrictions}
These are OS-9 commands under Unix names, and they behave as OS-9
commands. \cmd{mv} cannot move a file to another directory. \cmd{rmdir}
deletes the directory's contents too, after asking. \cmd{rm} takes no
options.

The Shell's syntax is OS-9's, not the Bourne shell's: the pipe symbol is
\texttt{!}, \texttt{>>} redirects standard error (there is no append),
\texttt{<>} and \texttt{>>>} redirect several paths at once, and
\texttt{\&} runs a command in the background.
\rh{Examples}
\begin{ex}
OS9:\textbf{ls -l /dd}
OS9:\textbf{mkdir /dd/src}
OS9:\textbf{cp notes /dd/src/notes}
OS9:\textbf{cd /dd/src}
OS9:\textbf{ps}
\end{ex}
\rh{Files}
\begin{keys}[1.5in]
\file{nitros9-l2/unix.tcl} & The list of aliases. Add a line to it and
rebuild to make another; see \rp{nitros9-l2}{8}.\\
\end{keys}
\rh{See Also} \rp{grep}{1}, \rp{pwd}{1}, \rp{vi}{1}
\end{refpage}

\refentry{vi}{1}{a small vi}
\begin{refpage}
\rh{Name} \cmd{vi}, \cmd{vim} -- a small vi
\rh{Syntax} \cmd{vi} [\texttt{-w} \var{cols}] [\texttt{-h} \var{rows}] [\var{file}]
\rh{Description}
The \cmd{vi} command is a screen editor in the manner of the Unix
\cmd{vi}, with most of the commands people use without thinking about
them. It starts in \emph{normal} mode, where keys are commands. The
insert commands enter \emph{insert} mode, where keys are text, until you
press \key{Esc}. A colon starts a \emph{command line} at the bottom of the
screen.

The bottom line shows messages, \texttt{-- INSERT --} in insert mode, and
the cursor's line and column. Lines longer than the screen are not wrapped:
the view scrolls sideways.

Most normal-mode commands take a count in front: \type{3dd} deletes three
lines, \type{5j} moves down five, \type{2d3w} deletes six words.

\sans{\textbf{Moving}}
\begin{keys}
\texttt{h j k l} & Left, down, up, right. The arrow keys too.\\
\texttt{w b e} & To the next word, back a word, to the end of the word.
\texttt{W B E} do the same with words that are anything but blanks.\\
\texttt{0 \^{} \$} & To the start of the line, its first non-blank, its end.\\
\texttt{G}, \texttt{gg} & To the last line, or to line \var{n} with a count;
to the first line.\\
\texttt{H M L} & To the top, middle, or bottom of the screen.\\
\texttt{f}\var{c} \texttt{F}\var{c} \texttt{t}\var{c} \texttt{T}\var{c} & To the next
(\texttt{f}) or previous (\texttt{F}) \var{c} on the line, or just before it
(\texttt{t}, \texttt{T}). \texttt{;} repeats; \texttt{,} repeats the other way.\\
\texttt{\%} & To the bracket that matches the next \texttt{( [ \{ \} ] )}.\\
\texttt{\{ \}} & Back or on a paragraph (to a blank line).\\
\texttt{+ - Return} & Down or up a line, to its first non-blank.\\
\texttt{\^{}F \^{}B} & Forward or back a screen. Page Down and Page Up too.\\
\texttt{\^{}D \^{}U} & Down or up half a screen.\\
\end{keys}

\sans{\textbf{Searching}}
\begin{keys}
\texttt{/}\var{text}, \texttt{?}\var{text} & Find \var{text} forward or backward,
wrapping around the file.\\
\texttt{n N} & Find it again, the same way or the other.\\
\texttt{*} & Find the word under the cursor.\\
\end{keys}

\sans{\textbf{Inserting}}
\begin{keys}
\texttt{i a} & Insert before, or append after, the cursor.\\
\texttt{I A} & Insert at the line's first non-blank, or append at its end.\\
\texttt{o O} & Open a new line below, or above.\\
\end{keys}
In insert mode, \key{Backspace} deletes the character before the cursor
(and can join lines), \texttt{\^{}W} the word before it, \texttt{\^{}U} the line
before it; \key{Return} splits the line, \key{Tab} inserts a tab, and the
arrow keys move the cursor. \key{Esc} goes back to normal mode.

\sans{\textbf{Changing}}
\begin{keys}
\texttt{x X} & Delete the character under, or before, the cursor.\\
\texttt{D C} & Delete, or change, to the end of the line.\\
\texttt{s S} & Change a character, or the whole line.\\
\texttt{r}\var{c} & Replace the character with \var{c} (\texttt{r} \key{Return} splits the line).\\
\texttt{J} & Join this line and the next.\\
\texttt{\~{}} & Switch the case of the character, and move on.\\
\texttt{>> <<} & Indent the line by a tab, or take one away.\\
\texttt{p P} & Put the last deleted or yanked text after, or before, the cursor
(below or above, for whole lines).\\
\texttt{u} & Undo the last change. Again and again, back to the start.\\
\texttt{.} & Do the last change again.\\
\end{keys}

\sans{\textbf{Operators}} \texttt{d} (delete), \texttt{c} (change),
\texttt{y} (yank), \texttt{>} and \texttt{<} (indent) work on the text that
any motion moves over: \type{dw}, \type{c\$}, \type{y3j}, \type{dG},
\type{d/end}, \type{>\}}. Doubled, they work on whole lines: \type{dd},
\type{cc}, \type{yy}, \type{>>}. \texttt{Y} is \texttt{yy}.

\sans{\textbf{Other}}
\begin{keys}
\texttt{ZZ} & Write the file if it was changed, and quit.\\
\texttt{ZQ} & Quit without writing.\\
\texttt{\^{}G} & Show the file's name, length, and where you are in it.\\
\texttt{\^{}L} & Redraw the screen.\\
\end{keys}

\sans{\textbf{Command line}}
\begin{keys}[1.75in]
\texttt{:w} [\var{file}] & Write the file, or write to \var{file}.\\
\texttt{:q}, \texttt{:q!} & Quit; \texttt{!} quits even if there are unsaved changes.\\
\texttt{:wq}, \texttt{:x} & Write and quit (\texttt{:x} writes only if there were changes).\\
\texttt{:e}[\texttt{!}] [\var{file}] & Edit another file, or the same file again.\\
\texttt{:r} \var{file} & Read \var{file} in below the current line.\\
\texttt{:}\var{n}, \texttt{:\$} & Go to line \var{n}, or the last line.\\
\texttt{:}[\var{range}]\texttt{s/}\var{old}\texttt{/}\var{new}\texttt{/}[\texttt{g}] &
Replace \var{old} with \var{new}, once in each line or everywhere
(\texttt{g}).\\
\texttt{:}[\var{range}]\texttt{d} & Delete lines.\\
\texttt{:help} & A reminder of these.\\
\end{keys}
A \var{range} is a line number, \texttt{.} (this line), \texttt{\$} (the
last), two of them with a comma (\texttt{2,10}), or \texttt{\%} (every
line). Without one, a command works on the current line.
\rh{Options}
\begin{keys}[0.8in]
\texttt{-w} \var{cols} & The terminal's width (default 80).\\
\texttt{-h} \var{rows} & The terminal's height (default 24).\\
\end{keys}
\rh{Restrictions}
Searches are for plain text, not regular expressions. There is one
register, there are no marks, no visual mode, no \texttt{:set}, and no
redo. The text must fit in memory: about 25K. Undo remembers as many
changes as fit in its own space, and forgets the oldest first.

OS-9 text files end their lines with a carriage return. \cmd{vi} reads
files that end them with CR, LF, or CR~LF, and writes CR.

The arrow keys work in insert mode because \cmd{vi} sets bit~7 of
\file{\$FF42} while it runs, so that the firmware sends them untranslated
(see \rp{tfr911h}{4}). If \cmd{vi} is killed, they stay that way until
the next \cmd{vi} ends.
\rh{Examples}
Change every \texttt{foo} to \texttt{bar}, write, and quit:
\begin{ex}
\textbf{:\%s/foo/bar/g}
\textbf{:wq}
\end{ex}
Delete from here to the line that contains \texttt{end}, then put it back:
\begin{ex}
\textbf{d/end}
\textbf{u}
\end{ex}
\rh{See Also} \rp{unix}{1}, \rp{tfr911h}{4}
\end{refpage}

% ─────────────────────────────────────────────────────────────────── 4
\refsection{4}{Special Files}
This section describes the hardware the TFR/911H firmware presents to NitrOS-9: the addresses in the 6309's I/O page that a program can read and write, and what they do.

\begin{tabularx}{\linewidth}{@{}>{\bfseries}p{1.3in}X@{}}
\toprule
\textnormal{\textbf{Name}} & \textbf{Description}\\
\midrule
fujinet(4) & The FujiNet: its wiring, the becker port, its disks, and booting from it\\
sn76489(4) & The sound chip\\
tfr911h(4) & The hardware the firmware emulates: memory, interrupts, console, disk, clock, keyboard\\
\bottomrule
\end{tabularx}


\refentry{fujinet}{4}{the FujiNet}
\begin{refpage}
\rh{Name} \cmd{fujinet} -- the FujiNet, its disks, and booting from it
\rh{Description}
A FujiNet is a network adapter for old computers: an ESP32 with WiFi and
an SD card. On the TFR/911H it runs the FujiNet firmware for the Color
Computer, built for the board \texttt{fujinet-coco-devkitc-centipede},
which speaks DriveWire over a 3.3-volt serial link at 921{,}600 baud. It is
wired as the Centipede's FujiNet hat is:

\begin{tabularx}{\linewidth}{@{}l l X@{}}
\toprule
\textbf{RP2350B} & \textbf{ESP32} & \textbf{Signal}\\
\midrule
GPIO~12 & GPIO~33 & Data to the FujiNet (UART0 TX to UART1 RX)\\
GPIO~13 & GPIO~21 & Data from the FujiNet (UART1 TX to UART0 RX)\\
GND & GND & \\
\bottomrule
\end{tabularx}

\sans{\textbf{The becker port}}\quad The firmware passes the link to the
6309 as a becker port, the kind CoCo emulators give DriveWire:

\begin{keys}[1.05in]
\file{\$FF45} & Read: bit~1 is set when a byte from the FujiNet is
waiting.\\
\file{\$FF46} & Read: the next byte from the FujiNet. Write: a byte to
the FujiNet.\\
\end{keys}
Emulators have the becker port at \file{\$FF41}, where the TFR/911H has its
sound chip, so NitrOS-9's DriveWire modules are built with
\texttt{BECKBASE} set to \file{\$FF45}. The firmware keeps 4K of what the
FujiNet sends, so nothing is lost while the 6309 is busy.

\sans{\textbf{Devices}}\quad \file{/X0}--\file{/X3} are the FujiNet's four
device slots, as DriveWire disks (\file{rbdw}). \file{/N} and
\file{/N1}--\file{/N5} are DriveWire's network devices (\file{scdwv}); the
\rp{fujinet}{1} commands use \file{/N1}, and its \texttt{SS.Fuji} SetStat,
which sends a request to the FujiNet and reads its reply with interrupts
masked. There is no virtual-serial poller, which would take the
FujiNet's replies for itself.

\sans{\textbf{Booting from the FujiNet}}\quad At the \texttt{TCL>} prompt,
\cmd{l2fn} mounts the image in the FujiNet's first device slot (slot~0),
and boots NitrOS-9 from it: the kernel is in the firmware, as for \cmd{l2},
and its \file{Boot} module loads \file{OS9Boot} from DriveWire drive~0.
That \file{OS9Boot} makes \file{/X0} the system disk, \file{/DD}; the
tether's disk is \file{/H0}. To make such a disk, put it on the FujiNet,
and boot it:
\begin{ex}
$ \textbf{cd v4/nitros9-l2}
$ \textbf{make fujinet-disk}
$ \textbf{curl -X MKCOL http://\var{fujinet}/dav/tfr911}
$ \textbf{curl -T level2-fn.dsk http://\var{fujinet}/dav/tfr911/level2-fn.dsk}
\end{ex}
Then, with NitrOS-9 booted with \cmd{l2}, put the image in device slot~0
(host slot~1 being the SD card), and boot it:
\begin{ex}
OS9: \textbf{fnsetdevfile 0 1 2 /tfr911/level2-fn.dsk}
\end{ex}
\begin{ex}
$ \textbf{BOOT=l2fn sh run-l2.sh}
\end{ex}
\cmd{make fujinet-disk} copies \file{level2.dsk}, with your files, to
\file{level2-fn.dsk}, and gives the copy the FujiNet's \file{OS9Boot}.
The FujiNet's web server offers its SD card by WebDAV at \file{/dav}.
\rh{Restrictions}
At power on, the FujiNet puts its own CONFIG disk in device slot~0; that
is why \cmd{l2fn} mounts the slot first.

The FujiNet keeps the image file open. After replacing the file, mount it
again (\cmd{l2fn} does, as does \texttt{fujinet mount 1} at
\texttt{TCL>}).

The FujiNet's firmware must read the version byte that NitrOS-9 sends after
\texttt{OP\_DWINIT}; firmware that does not loses the command after it, and
the 6309 waits for ever. See \file{v4/nitros9-l2/fujinet/README.md}.
\rh{Files}
\begin{keys}[2.3in]
\file{v4/tfr911h/v4\_fujinet.h} & The becker port and the link.\\
\file{v4/nitros9-l2/fujinet/} & The DriveWire modules, the \cmd{fn}
commands and their libraries, and how to set up the FujiNet.\\
\end{keys}
\rh{See Also} \rp{fnconfig}{1}, \rp{fujinet}{1}, \rp{tfr911h}{4}, \rp{tcl}{8}
\end{refpage}

\refentry{sn76489}{4}{the sound chip}
\begin{refpage}
\rh{Name} \cmd{sn76489} -- the sound chip
\rh{Syntax} \file{\$FF41}, write only
\rh{Description}
The TFR/911H firmware emulates a Texas Instruments SN76489 sound
generator at \file{\$FF41}, the address of the one in the Color Computer's
Game Master Cartridge. It has three square-wave tone voices and a noise
voice, each with sixteen volumes. The firmware mixes them at 31,250
samples a second, with the keyclick and bell sounds, and plays the result
as pulse-width-modulated audio on GPIO~8 (see Appendix~A).

The chip is programmed one byte at a time. A byte with bit~7 set
\emph{latches} a voice and a register and gives the low four bits of the
value; a byte with bit~7 clear gives the high six bits of a tone's period:

\begin{tabularx}{\linewidth}{@{}l X@{}}
\toprule
\textbf{Byte} & \textbf{Meaning}\\
\midrule
\texttt{1 vv 0 dddd} & Voice \texttt{vv} (0--2): low 4 bits of its period\\
\texttt{0 0 dddddd} & The high 6 bits of the latched voice's period\\
\texttt{1 vv 1 aaaa} & Voice \texttt{vv} (0--3): attenuation, 2~dB a step;
\texttt{1111} is off\\
\texttt{1 11 0 0 f rr} & The noise voice: \texttt{f}=1 white noise, 0
periodic; \texttt{rr} its rate, 0--2, or 3 to follow voice~2\\
\bottomrule
\end{tabularx}

A tone's period \var{N} (1--1023) gives a frequency of
3,579,545\,/\,(32\,×\,\var{N})~Hz, so A440 is \var{N}~=~254, and the
lowest note is about 109~Hz.

The control register at \file{\$FF42} is read and write:
\begin{keys}[0.6in]
bit 0 & Keyclicks on (see \rp{click}{1}).\\
bit 1 & The bell sounds on \texttt{\^{}G} (see \rp{beep}{1}).\\
bit 7 & Arrow keys untranslated (see \rp{tfr911h}{4}).\\
\end{keys}
\rh{Examples}
In C, compiled with CMOC, from \rp{tetris}{6}:
\begin{ex}
#define SND (*(volatile unsigned char *)0xFF41)

void note(unsigned char voice, unsigned n) \{   /* n = 3579545/32/Hz */
    SND = 0x80 | (voice << 5) | (n & 15);
    SND = (n >> 4) & 0x3F;
\}

void volume(unsigned char voice, unsigned char att) \{ /* 15 is off */
    SND = 0x90 | (voice << 5) | att;
\}
\end{ex}
\rh{See Also} \rp{beep}{1}, \rp{click}{1}, \rp{tfr911h}{4}, Appendix~A
\end{refpage}

\refentry{tfr911h}{4}{the hardware the firmware emulates}
\begin{refpage}
\rh{Name} \cmd{tfr911h} -- the hardware the TFR/911H firmware emulates for NitrOS-9 Level 2
\rh{Description}
To NitrOS-9, the TFR/911H is a small Color Computer~3. The firmware's
second core answers every bus cycle of the 6309: reads and writes of RAM
go to the RP2350B's memory through the emulated MMU, and those of the I/O
page, \file{\$FF00}--\file{\$FFFF}, go to the devices below. Locations not
listed read back what was last written to them. The 6309 runs in 6809
mode.

\sans{\textbf{Memory}}\quad 256K of RAM, in 32 blocks of 8K. The GIME's MMU
maps eight blocks into each of two tasks' 64K address spaces:
\file{\$FFA0}--\file{\$FFA7} for task~0 and \file{\$FFA8}--\file{\$FFAF}
for task~1. \file{\$FF90} bit~6 enables the MMU and bit~5 the GIME's
interrupts; \file{\$FF91} bit~0 picks the task. \file{\$FE00}--\file{\$FEFF}
is the constant page, the same in every map.

\sans{\textbf{Interrupts}}\quad \file{\$FF92} bit~3 enables the vertical
border (VBORD) interrupt, which a timer in the firmware raises 60 times a
second; reading \file{\$FF92} acknowledges it. The 6850's interrupt joins
it on the 6309's IRQ line. There is no FIRQ.

\sans{\textbf{Console}}\quad An MC6850 ACIA at \file{\$FF14} (status and
control) and \file{\$FF15} (data). Its transmitter is always ready;
characters go to the tether at once.

\sans{\textbf{Disk}}\quad EmuDsk, a virtual hard disk whose sectors the
tether reads and writes in a file:
\begin{keys}[1.05in]
\file{\$FF80}--\file{\$FF82} & The logical sector number, 24 bits.\\
\file{\$FF83} & Write a command: 0 reads a sector, 1 writes one. Read the
status: 0 if it worked.\\
\file{\$FF84}--\file{\$FF85} & The address of the 256-byte buffer.\\
\file{\$FF86} & The drive number.\\
\end{keys}
The bus waits while the tether moves the sector over USB, so the status is
ready in the instruction after the command.

\sans{\textbf{Clock}}\quad \file{\$FFC0}--\file{\$FFC7} hold the century,
year, month, day, day of the week, hour, minute, and second, in binary,
in the format of Jeff Vavasour's CoCo~3 emulator. The tether sends your
computer's local time when it connects, and the firmware keeps it going.

\sans{\textbf{Sound}}\quad \file{\$FF41} and \file{\$FF42}; see \rp{sn76489}{4}.

\sans{\textbf{FujiNet}}\quad A becker port at \file{\$FF45} (status) and
\file{\$FF46} (data), for DriveWire to a FujiNet; see \rp{fujinet}{4}.

\sans{\textbf{Info port}}\quad The firmware's report on itself, for
\rp{info}{1}. Write \file{\$A5} to \file{\$FF43}; when \file{\$FF43} reads
1, each read of \file{\$FF44} gives the report's next character, and 0
after the last. The report is lines of the form
\var{label}~\texttt{TAB}~\var{value}~\texttt{LF}.

\sans{\textbf{Keyboard}}\quad Keys typed in the tether reach OS-9 through the
6850. The tether turns Delete into Backspace and Enter into a carriage
return. The arrow keys arrive as a Color Computer~3 sends them, unless
bit~7 of \file{\$FF42} is set:

\begin{tabularx}{\linewidth}{@{}l l l X@{}}
\toprule
\textbf{Key} & \textbf{CoCo code} & \textbf{Bit 7 set} & \textbf{Note}\\
\midrule
Up & \file{\$0C} & 128 & \\
Down & \file{\$0A} & 129 & \\
Left & \file{\$08} & 130 & The same as Backspace, as CoCo code\\
Right & \file{\$09} & 131 & The same as Tab, as CoCo code\\
Page Up & --- & 132 & Dropped, as CoCo code\\
Page Down & --- & 133 & Dropped, as CoCo code\\
\bottomrule
\end{tabularx}
\rh{Files}
\begin{keys}[2.2in]
\file{v4/tfr911h/v4\_nitros9\_l2.h} & The emulation.\\
\file{v4/tfr911h/v4\_audio.h} & The sound chip and the mixer.\\
\file{v4/tfr911h/v4\_fujinet.h} & The becker port.\\
\end{keys}
\rh{See Also} \rp{info}{1}, \rp{fujinet}{4}, \rp{sn76489}{4}, \rp{tether}{8}
\end{refpage}

% ─────────────────────────────────────────────────────────────────── 6
\refsection{6}{Games}
This section describes the games. Both draw with ANSI escape sequences on the terminal the tether runs in.

\begin{tabularx}{\linewidth}{@{}>{\bfseries}p{1.3in}X@{}}
\toprule
\textnormal{\textbf{Name}} & \textbf{Description}\\
\midrule
tetris(6) & Tetris, with music\\
zm(6) & A Z-machine, for Zork and the other Infocom games\\
\bottomrule
\end{tabularx}


\refentry{tetris}{6}{Tetris, with music}
\begin{refpage}
\rh{Name} \cmd{tetris} -- Tetris, with music
\rh{Syntax} \cmd{tetris}
\rh{Description}
Pieces fall into a well ten cells wide and twenty deep. Move and turn them
so they make full rows, which disappear. The game draws in color on the
terminal the tether runs in, shows the next piece and a ghost of where the
falling one will land, and deals the seven pieces from a shuffled bag, so
no piece is ever long in coming.
\begin{keys}[1.3in]
\key{Left} \key{Right}, \texttt{a d} & Move the piece.\\
\key{Up}, \texttt{w} & Turn it.\\
\key{Down}, \texttt{s} & Drop it a row.\\
\key{Space} & Drop it all the way.\\
\texttt{p} & Pause.\\
\texttt{m} & Turn the music off or on.\\
\texttt{q} & Quit.\\
\end{keys}
Clearing one, two, three, or four rows at once scores 40, 100, 300, or
1,200 points times the level plus one. Dropping a piece scores a point a
row, or two with \key{Space}. Every ten rows is a new level, and the
pieces fall faster, up to level~19.

With a speaker connected, the game plays \emph{Korobeiniki} on the sound
chip, faster as the levels go up, with sounds for moving, turning,
landing, clearing rows, a new level, and the end of the game.
\rh{Restrictions}
It needs an ANSI color terminal of at least 80 by 24. It turns keyclicks
off while it runs.
\rh{See Also} \rp{click}{1}, \rp{sn76489}{4}
\end{refpage}

\refentry{zm}{6}{a Z-machine, for Infocom's games}
\begin{refpage}
\rh{Name} \cmd{zm} -- play Infocom's games, and others written for the Z-machine
\rh{Syntax} \cmd{zm} [\texttt{-w} \var{cols}] [\texttt{-h} \var{rows}] [\texttt{-m} \var{kbytes}] [\texttt{-d}] \var{story}
\rh{Description}
Infocom wrote its games for an imaginary computer, the Z-machine, and sold
them with an interpreter for each real one. \cmd{zm} is a Z-machine
interpreter for OS-9. It plays \emph{version~3} story files: most of
Infocom's games from 1980 to 1987 (Zork I--III, The Hitchhiker's Guide to
the Galaxy, Planetfall, Enchanter, Deadline, and many more), and games
written with Inform and compiled with \texttt{-v3}.

\cmd{zm} looks for \var{story}, then \var{story}\texttt{.z3}, and, if the
name has no directory in it, the same in \file{/DD/GAMES}. So \type{zm
zork1} finds \file{/DD/GAMES/zork1.z3}.

The top line of the screen is the status line: where you are, and your
score and moves, or the time in the games that keep it. The text scrolls
beneath it, word-wrapped to the screen's width, and pauses with
\texttt{[MORE]} when a screen goes by without your typing. Games that
split the screen, as Seastalker does, can.

The game's commands \texttt{save} and \texttt{restore} ask for a file name,
offering the story's name with \texttt{.sav}; save files go in the data
directory. \texttt{restart}, \texttt{verify}, and \texttt{quit} work as in
the games' manuals.

The part of the story the game can change, its \emph{dynamic memory}, is
kept in RAM. So is as much of the rest as fits---about 40K, got by growing
the process's data area with \texttt{F\$Mem}. If the story is bigger than
that, the rest is read from the file in 512-byte pages as the game needs
them, and the least recently used page gives way. Panic, which is 37K, is
held in memory whole; Zork~I, at 87K, is paged, and still answers a
command in under a second.
\rh{Options}
\begin{keys}[0.8in]
\texttt{-w} \var{cols} & The terminal's width (default 80).\\
\texttt{-h} \var{rows} & The terminal's height (default 24).\\
\texttt{-m} \var{kbytes} & Use at most this much memory for the story
(8--60).\\
\texttt{-d} & A dumb terminal: no status line, no escape sequences, no
\texttt{[MORE]}.\\
\end{keys}
\rh{Restrictions}
Only version~3 stories are played; \cmd{zm} says which version a story is
if it is another. There is no transcript (\texttt{script}), and there are
no sounds but the two beeps of The Lurking Horror, which it plays on the
sound chip.
\rh{Diagnostics}
If the story does something impossible, \cmd{zm} stops with a message
such as:
\begin{ex}
*** zm: stack underflow (PC $1A2B3) ***
\end{ex}
\rh{Examples}
\begin{ex}
OS9:\textbf{zm zork1}
ZORK I: The Great Underground Empire
...
West of House
You are standing in an open field west of a white house, with a boarded
front door.
There is a small mailbox here.

>\textbf{open mailbox}
\end{ex}
To put story files on the disk, see \rp{nitros9-l2}{8}.
\rh{Files}
\begin{keys}[1.6in]
\file{/DD/GAMES/*.z3} & Story files.\\
\var{story}\file{.sav} & Saved games.\\
\end{keys}
\rh{See Also} \rp{nitros9-l2}{8}, Appendix~B
\end{refpage}

% ─────────────────────────────────────────────────────────────────── 8
\refsection{8}{Maintenance}
This section describes the programs that run on the computer the TFR/911H is plugged into, the firmware's console, and how to build the system and keep a disk up to date.

\begin{tabularx}{\linewidth}{@{}>{\bfseries}p{1.3in}X@{}}
\toprule
\textnormal{\textbf{Name}} & \textbf{Description}\\
\midrule
nitros9-l2(8) & Build NitrOS-9 Level 2 for the TFR/911H, and keep a disk up to date\\
tcl(8) & The firmware's console\\
tether(8) & The board's terminal, disk, and clock\\
\bottomrule
\end{tabularx}


\refentry{nitros9-l2}{8}{build the system and keep a disk up to date}
\begin{refpage}
\rh{Name} \cmd{nitros9-l2} -- build NitrOS-9 Level 2 for the TFR/911H, and keep a disk up to date
\rh{Syntax} \cmd{make} [\var{target}] [\texttt{DISK=}\var{image}] [\texttt{STORIES=}\var{files}] [\texttt{OS9\_TOOL=}\var{path}]
\rh{Description}
The \file{v4/nitros9-l2} directory builds NitrOS-9 Level 2 for the TFR/911H
from a NitrOS-9 source tree, with the commands described in this manual,
and makes a disk image, and another that boots from a FujiNet. It needs \cmd{lwasm} from LWTOOLS, \cmd{os9} from
ToolShed, \cmd{tclsh}, \cmd{python3}, and, for the C programs, CMOC.

\begin{keys}[1.35in]
\texttt{make} & Build the kernel, \file{OS9Boot}, the commands, and
\file{build/tfr911-level2.dsk}, and the firmware's copy of the kernel
track, \file{../tfr911h/nitros9\_l2\_rom.h}.\\
\texttt{make install} & Bring the disk \var{DISK} (default
\file{level2.dsk}) up to date: a new \file{OS9Boot}, and every file in
\file{CMDS} and \file{SYS}. Other files are left alone.\\
\texttt{make stories} & Copy the story files named by \var{STORIES} to
\file{/DD/GAMES} on \var{DISK}.\\
\texttt{make fujinet-disk} & Copy \var{DISK} to \var{FNDISK} (default
\file{level2-fn.dsk}), with the \file{OS9Boot} that boots from the
FujiNet; see \rp{fujinet}{4}.\\
\texttt{make build/zm-host} & Build \cmd{zm} for this computer, for
testing; also \texttt{build/vi-host}.\\
\texttt{make clean} & Remove the build.\\
\end{keys}

The firmware is built from \file{v4}: \type{make build/tfr911h\_v4.uf2},
then loaded with \cmd{picotool}.
\rh{Restrictions}
Quit the tether before changing the disk it is serving, and reset the
board after \texttt{make install} so that it boots the new
\file{OS9Boot}.
\rh{Examples}
\begin{ex}
$ \textbf{cd v4/nitros9-l2}
$ \textbf{make OS9_TOOL=~/toolshed/build/unix/os9/os9}
$ \textbf{make install DISK=level2.dsk}
$ \textbf{make stories STORIES="~/if/zork1.z3 ~/if/hhgg.z3" DISK=level2.dsk}
\end{ex}
\rh{Files}
\begin{keys}[1.45in]
\file{rmake-l2.tcl} & Which modules go in \file{OS9Boot} and on the disk.\\
\file{unix.tcl} & The Unix-style aliases, and changes to the Shell and \cmd{grep}.\\
\file{modules/} & The TFR/911H's own modules: the clock, \cmd{clear}, \cmd{beep}, \cmd{click}.\\
\file{fujinet/} & The DriveWire modules and the \cmd{fn} commands.\\
\file{games/} & The C programs: \cmd{tetris}, \cmd{zm}, \cmd{vi}, \cmd{info}.\\
\file{helpmsg.unix} & The \cmd{man} pages for what was added.\\
\file{doc/manual.tex} & This manual.\\
\end{keys}
\rh{See Also} \rp{tcl}{8}, \rp{tether}{8}
\end{refpage}

\refentry{tcl}{8}{the firmware's console}
\begin{refpage}
\rh{Name} \cmd{tcl} -- the TFR/911H firmware's console
\rh{Syntax} \texttt{TCL>} \var{command}
\rh{Description}
When the TFR/911H starts, its firmware waits at the \texttt{TCL>} prompt,
a Tcl interpreter that runs on the RP2350B. Besides the Tcl language, it
has these commands:

\begin{keys}[1.3in]
\cmd{l2} & Boot NitrOS-9 Level 2 from EmuDsk drive~0.\\
\cmd{l2fn} [\texttt{-n}] & Mount the image in the FujiNet's device slot~0
and boot NitrOS-9 Level 2 from it; with \texttt{-n}, boot from whatever
is mounted. See \rp{fujinet}{4}.\\
\cmd{fujinet} & Ask the FujiNet the time, to show that the link works.\\
\cmd{fujinet mount} [\var{n}] & Mount the image in the FujiNet's device
slot \var{n} (1--4, the FujiNet's own numbers; default 1).\\
\cmd{fujinet stats} & The link's byte and error counts.\\
\cmd{bye} & Boot TurbOS, the OS-9 Level 1 system in the firmware.\\
\cmd{beep} [\var{hz} [\var{ms}]] & Play a tone through the speaker
(default 880~Hz for 200~ms), to test the wiring.\\
\cmd{click} [\texttt{on|off}] & Turn keyclicks on or off, and say which.\\
\end{keys}
\rh{Examples}
\begin{ex}
TCL> \textbf{beep 440 1000}
TCL> \textbf{fujinet}
FujiNet says it is 2026-10-17 18:46:20
TCL> \textbf{l2}
\end{ex}
\rh{See Also} \rp{click}{1}, \rp{fujinet}{4}, \rp{tether}{8}
\end{refpage}

\refentry{tether}{8}{the board's terminal, disk, and clock}
\begin{refpage}
\rh{Name} \cmd{tether}, \cmd{run-l2.sh} -- the board's terminal, disk, and clock
\rh{Syntax} \cmd{sh run-l2.sh}\\
\cmd{tether} \texttt{-wire} \var{device} [\texttt{-disks} \var{images}] [\texttt{-boot} \var{command}] [\texttt{-pass\_ctrl\_c}] [\var{other options}]
\rh{Description}
The tether is a program on your computer that talks to the TFR/911H over
USB. It is the board's terminal: what you type goes to the 6309, and what
the 6309 prints appears in the window. It is the board's disk: it reads and
writes the sectors of disk image files for EmuDsk. And it sets the board's
clock when it connects.

\file{v4/run-l2.sh} builds the tether and runs it with the options
NitrOS-9 Level 2 needs. If NitrOS-9 is already running on the board, it
simply reconnects.

\begin{keys}[1.25in]
\texttt{WIRE} & The board's USB serial device. On macOS, the only
\file{/dev/cu.usbmodem*} device, if there is just one.\\
\texttt{DISK} & The disk image for drive~0. The default,
\file{nitros9-l2/level2.dsk}, is copied from the build the first time.\\
\texttt{BOOT} & What to type at \texttt{TCL>}: \texttt{l2} (the default),
or \texttt{l2fn} to boot from the FujiNet.\\
\end{keys}
\rh{Options}
\begin{keys}[1.25in]
\texttt{-wire} \var{device} & The serial device.\\
\texttt{-disks} \var{images} & Disk image files, comma separated, for
drives 0, 1, \ldots\\
\texttt{-boot} \var{command} & A command to type at the \texttt{TCL>}
prompt, such as \texttt{l2}. Ignored if the 6309 is already running.\\
\texttt{-pass\_ctrl\_c} & Send \texttt{\^{}C} (and \texttt{\^{}\textbackslash} and \texttt{\^{}Z}) to
the 6309, and quit the tether with \texttt{\^{}]} instead.\\
\texttt{-level} \var{n} & The NitrOS-9 level, for the tether's
diagnostics.\\
\end{keys}
\rh{Restrictions}
Only one program can have the board's serial device open. Quit the tether
before you flash the firmware or change a disk image it is serving.
\rh{Files}
\begin{keys}[1.25in]
\file{v4/\_log} & The tether's diagnostics, including the firmware's
statistics every ten seconds.\\
\end{keys}
\rh{See Also} \rp{nitros9-l2}{8}, \rp{tcl}{8}, \rp{tfr911h}{4}
\end{refpage}

% ═══════════════════════════════════════════════════════════════════
\clearpage
\appendix
\renewcommand{\therefname}{}
\chapter{Connecting a Speaker}
\chapterpages
\renewcommand{\thesectname}{Connecting a Speaker}

The TFR/911H makes sound as pulse-width-modulated audio on GPIO~8 of the
RP2350B: an 8-bit level carried at about 977~kHz, far above hearing. A
filter and a small amplifier turn it into sound. The firmware was written
for an Adafruit PAM8302, a 2.5~W class~D amplifier on a small board.

\section{Wiring}

GPIO~8 is pin~1 of the 2-by-10 header J3, and is not connected to the
6309. Wire the PAM8302 like this:

\begin{tabularx}{\linewidth}{@{}l l X@{}}
\toprule
\textbf{J3 pin} & \textbf{Signal} & \textbf{PAM8302}\\
\midrule
2 & VBUS (5~V) & Vin\\
6 & GND & GND, and A−\\
1 & GPIO 8 & A+, through a 1~kΩ resistor, with a 10~nF capacitor from A+ to GND\\
\bottomrule
\end{tabularx}

The resistor and capacitor make a low-pass filter that takes out the
carrier. Connect the speaker to the PAM8302's + and − terminals, and
turn its volume control down before you start.

\section{Trying It}

At the \texttt{TCL>} prompt, before an OS is booted:
\begin{ex}
TCL> \textbf{beep 440 1000}
\end{ex}
should play an A for a second. Then, in NitrOS-9, \cmd{beep} rings the
bell, every key clicks (\cmd{click off} stops that), and \cmd{tetris}
plays music.

\section{How It Works}

A timer interrupt on the RP2350B's first core runs 31,250 times a second.
Each time, it takes any bytes the 6309 has written to the SN76489, works
out the next sample of the three tone voices, the noise voice, and the
keyclick or bell, adds them, and sets the PWM level. The 6309's writes
reach it through a queue, so the core that runs the bus never waits.

% ═══════════════════════════════════════════════════════════════════
\chapter{How the Software Was Tested}
\chapterpages
\renewcommand{\thesectname}{How the Software Was Tested}

The C programs were written so that the same source builds for OS-9 with
CMOC and for a Unix computer with an ordinary C compiler. That let each one
be tested against a reference on the big computer, and then on the board.

\section{zm}

\begin{itemize}\setlength{\itemsep}{0.3ex}
\item CZECH, the Comprehensive Z-machine Emulation Checker, compiled for
version~3 with Inform: all 349 of its tests pass, on the Unix build and on
the TFR/911H's 6309, where the run takes about five seconds. They pass
too with memory cut to 8K, which makes nearly every read a page from the
file.
\item Zork I, II, and III and Inform's Adventure: transcripts of
the same commands give the same text, word for word, as \cmd{dfrotz}, a
widely used interpreter, with the story wholly in memory and with only
16K to 40K for it.
\item On the board: Zork~I answers each command in 0.4 to 0.9 seconds;
\texttt{save}, \texttt{restore}, \texttt{restart} work.
\end{itemize}

\section{vi}

\begin{itemize}\setlength{\itemsep}{0.3ex}
\item 123 editing commands---motions, operators, counts, undo,
\texttt{.}, searches, \texttt{:s} and ranges---were typed into Vim and
into \cmd{vi}, on a pseudo-terminal at the same pace, and the files they
wrote compared: all the same.
\item 27 sessions on a 300-line file with tabs and long lines checked that
what \cmd{vi} shows is what is in the file, after scrolling, sideways
scrolling, and edits.
\item On the board: moving and typing answer in under 0.05 seconds,
scrolling in 0.15, and replacing 178 strings on 116 lines in 1.8.
\end{itemize}

\section{info}

Its report was compared with what NitrOS-9's own \cmd{procs}, \cmd{mdir},
\cmd{mfree}, and \cmd{free} say, on the board, and with a program running
in the background.

\section{The Firmware}

Each change was tried on the board: NitrOS-9 boots, the disk and clock
work, Tetris's writes reach the sound chip (182 of them in one game, by
the firmware's count), and the arrow keys arrive as they should with bit~7
of \file{\$FF42} clear and set.

\end{document}
